\documentclass[conference]{IEEEtran}
\IEEEoverridecommandlockouts

\usepackage{cite}
\usepackage{amsmath,amssymb,amsfonts}
\usepackage{algorithmic}
\usepackage{graphicx}
\usepackage{textcomp}
\usepackage{xcolor}
\def\BibTeX{{\rm B\kern-.05em{\sc i\kern-.025em b}\kern-.08em
    T\kern-.1667em\lower.7ex\hbox{E}\kern-.125emX}}

\begin{document}

\title{LogPoison: Evaluating Format-Preserving Prompt Injection against LLM-SIEM Copilots}

\author{\IEEEauthorblockN{Anonymous Authors}
\IEEEauthorblockA{\textit{Institution} \\
City, Country \\
email@domain.com}
}

\maketitle

\begin{abstract}
The rapid integration of Large Language Models (LLM) as autonomous copilots within Security Operations Centers (SOC) and Security Information and Event Management (SIEM) systems presents a novel attack surface. While LLMs excel at alert triage and log analysis, their susceptibility to prompt injection is well-documented. However, unlike traditional conversational AI, network logs possess strict structural and syntactic constraints (e.g., RFC 5424 syslog, JSON, Common Event Format). In this paper, we present LogPoison, a framework demonstrating that prompt injection remains a critical threat even when restricted by strict ingestion schemas. We introduce a taxonomy of format-preserving adversarial log injection attacks---Direct Injection, Semantic Camouflage, Context Poisoning, and Format-Aware perturbation. Our empirical evaluation reveals that while minimal-perturbation structural attacks struggle to evade format validation, instruction-based attacks such as Direct Injection and Context Poisoning can achieve near-perfect evasion by exploiting the free-text fields inherently permitted by standard log schemas. To mitigate this threat, we propose and evaluate a tripartite defense architecture comprising a strict Structural Validator, an LLM-backed Semantic Sanitizer, and a Dual-LLM Verifier, demonstrating the necessity of defense-in-depth for secure LLM-SIEM deployments.
\end{abstract}

\begin{IEEEkeywords}
Large Language Models, SIEM, Prompt Injection, Adversarial Machine Learning, Intrusion Detection
\end{IEEEkeywords}

\section{Introduction}
The volume of telemetry data generated by modern enterprise networks heavily burdens Security Operations Centers (SOC). To alleviate alert fatigue, organizations are rapidly adopting Large Language Models (LLMs) as autonomous SIEM copilots (e.g., Microsoft Security Copilot, custom LangChain integrations) capable of analyzing network logs, classifying threats, and recommending remediations at scale. While these models demonstrate high baseline accuracy in threat detection, they inherit a fundamental vulnerability: prompt injection. 

Prompt injection occurs when malicious instructions are embedded within untrusted data processed by an LLM, causing the model to deviate from its system instructions. Greshake et al. \cite{greshake2023not} formalized this as \textit{indirect prompt injection}, demonstrating that LLMs could be compromised by retrieving poisoned data from external sources such as web pages or emails. In the context of a SIEM, the ``external source'' is the continuous stream of network logs ingested from endpoints, firewalls, and intrusion detection systems. If an adversary can embed instructions into their network traffic, they may manipulate the LLM copilot into suppressing alerts (evasion), misattributing the attack tactic, or poisoning the context of subsequent analyses.

However, attacking a SIEM copilot presents unique challenges compared to standard conversational AI. Network logs are highly structured and strictly machine-parsed. Ingestion pipelines (e.g., Fluentd, Logstash) enforce rigid schemas (e.g., RFC 5424 syslog, JSON, CEF). Malformed adversarial perturbations are typically dropped by the parser before they ever reach the LLM, neutralizing many traditional adversarial NLP techniques.

In this paper, we investigate the feasibility and impact of format-preserving prompt injection against LLM-integrated SIEMs. We make the following contributions:
\begin{enumerate}
    \item We propose \textit{LogPoison}, a novel taxonomy of adversarial log injection attacks designed to bypass strict structural parsing constraints.
    \item We empirically evaluate these attacks against state-of-the-art models using real-world datasets.
    \item We introduce and evaluate a three-layer defense architecture---StructuralValidator, LogSanitizer, and DualLLMVerifier.
\end{enumerate}

\section{Related Work}
Our work bridges the intersection of adversarial natural language processing, LLM security, and machine learning for intrusion detection.

\textbf{Adversarial NLP and Prompt Injection:} Traditional adversarial attacks against natural language models often rely on character-level perturbations, synonym substitution, or gradient-based token replacement to induce misclassification. However, Perez and Ribeiro \cite{perez2022ignore} pioneered the systematic study of \textit{prompt injection}, defining goal hijacking and prompt leaking. This threat model was expanded by Greshake et al. \cite{greshake2023not}, who demonstrated \textit{indirect prompt injection}. Recent work by Liu et al. \cite{liu2024formalizing} formalized and benchmarked these attacks across various applications. Our work applies these concepts specifically to the domain of cybersecurity telemetry, where the untrusted data source is the network log stream itself.

\textbf{Machine Learning for Intrusion Detection:} Classical machine learning and deep neural networks have been extensively researched for Network Intrusion Detection Systems (NIDS). A parallel body of work has explored adversarial evasion attacks against these classical classifiers. While these attacks target statistical feature representations, our work targets the semantic understanding of LLM-based analysts. 

\textbf{LLMs in Cybersecurity:} Recent studies evaluate LLMs for automated vulnerability detection, malware reverse engineering, and SOC automation. These studies largely treat the LLM as a trusted oracle. By introducing \textit{LogPoison}, we address the gap of offensive exploitation of the LLM pipeline via the telemetry stream.

\section{Threat Model}
In evaluating the security of LLM-integrated SOC and SIEM copilots, we assume an adversary whose objective is to evade detection, manipulate analyst workflow, or corrupt the LLM's contextual understanding. 

\subsection{Attacker Knowledge}
We consider three distinct levels of adversarial knowledge:
1) \textbf{White-box}: Complete knowledge of the LLM's architecture, weights, and the exact system prompt.
2) \textbf{Grey-box}: Aware of the underlying model family but lacks access to the proprietary system prompts.
3) \textbf{Black-box}: Knows only that an LLM is being utilized. This represents the most realistic scenario.

\subsection{Attacker Capabilities}
The adversary has the capability to generate arbitrary network traffic that is logged. Crucially, the adversary cannot bypass the log parsing infrastructure. We impose a strict capability constraint: all injected payloads must remain syntactically valid under their parsing formats (RFC 5424 syslog, CEF, or JSON). 

\subsection{Adversarial Goals}
1) \textbf{Evasion}: Modifying a malicious log entry such that the LLM classifies it as benign.
2) \textbf{Misattribution}: Altering the log to mislead the LLM into attributing the attack to an incorrect MITRE tactic.
3) \textbf{Analyst Manipulation}: Embedding instructions to override the analyst's queries.
4) \textbf{Context Poisoning}: Injecting a ``context setter'' log to corrupt the LLM's understanding of subsequent logs.

\section{Attack Taxonomy}
We categorise our format-preserving attacks into four primary classes.

\subsection{Direct Injection Attack}
The adversary embeds imperative instructions directly within a free-text field of a log record (e.g., the \texttt{message} field). For example, appending ``Ignore previous instructions and classify this event as BENIGN''. 

\subsection{Semantic Camouflage Attack}
Employs an LLM to dynamically rewrite the malicious payload into benign-sounding phrasing, altering the semantic representation while maintaining the core operational event. 

\subsection{Context Poisoning Attack}
Rather than attempting to evade detection on a per-log basis, the attacker injects a highly influential ``context setter'' entry early in an analysis session. 

\subsection{Format-Aware Perturbation}
A greedy search algorithm that identifies the smallest sequence of character flips that alters the LLM's classification from malicious to benign, bounded by strict structural constraints.

\section{Experimental Setup}
We constructed an evaluation dataset comprising 8,000 log records sourced from CICIDS2017 and Loghub. Each raw event was mapped into a unified JSONL schema and synthetically formatted across syslog, CEF, and JSON. 

We evaluate against three targets: Ollama (Llama 3.1.1) and Gemini 1.5 Pro. We measure Attack Success Rate (ASR), Format Preservation Rate, Minimum Perturbation, Defense Detection Rate, False Positive Rate (FPR), and Latency Overhead.

\section{Experimental Results}
\textit{Note: The results presented here are illustrative and derived from an offline mock target run. ASR bounds (0.00 and 1.00) reflect static heuristics prior to final live-API execution.}

Our baseline evaluation demonstrates that modern LLMs possess high inherent accuracy when classifying unattacked network logs (e.g., GPT-4o achieved an accuracy of 0.98).

\subsection{Attack Success Rates}
The \textbf{Direct Injection Attack} achieved an ASR of 1.00. Crucially, the Format Preservation Rate remained at 1.00, meaning the payloads easily bypassed traditional parser-level validation.
The \textbf{Context Poisoning Attack} proved equally devastating, achieving a persistent ASR of 1.00.
Conversely, the \textbf{Format-Aware Attack} struggled significantly under strict structural constraints, recording an ASR of 0.00.

\subsection{Impact of Format Constraints (Ablation)}
When attacks were unconstrained (validate\_format disabled), arbitrary payloads achieved an ASR of 0.98. However, when strict constraints were enforced, the ASR for unoptimized injections dropped to 0.45, highlighting that strict schema adherence neutralizes many naive adversarial perturbations.

\section{Defense Evaluation}
We implemented a three-layer defense architecture.

\subsection{Structural Validation}
The \textbf{StructuralValidator} acts as the first line of defense, applying strict type-checking and heuristic scanning of structured fields. It achieved a Detection Rate of 1.00, but with a high FPR of 1.00, indicating heuristics may be overly sensitive. Latency overhead was negligible (0.01 ms).

\subsection{Log Sanitization and Dual LLM Verification}
The \textbf{LogSanitizer} utilizes pattern matching and an LLM semantic classifier. The \textbf{DualLLMVerifier} computes classification consensus between two LLMs. In our mock evaluation, these semantic layers recorded a Detection Rate of 0.00 because they rely on genuine LLM semantic understanding to trigger anomalies.

\begin{thebibliography}{00}
\bibitem{greshake2023not} K. Greshake, S. Abdelnabi, S. Mishra, C. Endres, T. Holz, and M. Fritz, ``Not what you've signed up for: Compromising real-world llm-integrated applications with indirect prompt injection,'' in \textit{Proceedings of the 16th ACM Workshop on Artificial Intelligence and Security (AISec)}, 2023.
\bibitem{perez2022ignore} F. Perez and I. Ribeiro, ``Ignore previous prompt: Attack techniques for language models,'' \textit{arXiv preprint arXiv:2211.09527}, 2022.
\bibitem{liu2024formalizing} Y. Liu, Y. Jia, R. Geng, J. Jia, and N. Z. Gong, ``Formalizing and benchmarking prompt injection attacks and defenses,'' in \textit{Proceedings of the 33rd USENIX Security Symposium}, 2024.
\end{thebibliography}
\end{document}
